\section{Practical Implications for UHS Surrogate Modeling}
The development and evaluation of GNN-based emulators provide several practical lessons for utilizing Scientific Machine Learning in Underground Hydrogen Storage operations.

\subsection{Current Capabilities and screening Studies}
The principal advantage of the GNN-based surrogate is its computational efficiency. While a full numerical reservoir simulation of Case 5 requires minutes to hours of wall-clock time, the trained ST-GNN v2 performs a 142-step autoregressive rollout in milliseconds. This substantial acceleration makes the surrogate highly suitable for rapid screening studies. For instance, reservoir engineers can evaluate hundreds of cyclic injection-withdrawal schedule permutations to maximize hydrogen recovery factor or identify optimal operating envelopes, under the constraint that the geostatistical realization and well locations remain fixed. The model captures the general pressure dynamics and phase-change transitions, identifying whether a schedule will lead to excessive pressure buildup or depletion.

\subsection{Unsuitability for Direct Field Control}
Despite its speed, the current surrogate is unsuitable for direct field application or safety-critical decisions. In UHS operations, pressure management is critical to prevent caprock fracturing, fault reactivation, and hydrogen leakage. The ST-GNN v2 model exhibits a peak pressure RMSE of 59.19~bar at step 50, which corresponds to a 49.0\% deviation relative to the average reservoir operating pressure of 120.70~bar. An error of this magnitude is unacceptable for geomechanical integrity monitoring or real-time pressure control, where pressure constraints must be maintained within tight safety margins.

Furthermore, the surrogate co-emulates gas saturation ($S_g$) and pressure change, but does not enforce physical conservation laws. During long rollouts, accumulated saturation errors (reaching a peak RMSE of 0.404) lead to predictions that violate mass balance. For direct field control, emulators must provide guaranteed physical bounds, which requires integrating hard physical constraints or mass-conservation layers into the network.

\subsection{Dataset Expansion and Generalization Challenges}
A critical requirement for moving beyond localized cyclic emulation is expanding the training dataset. The current model is trained on only four simulation cases and evaluated on a single geostatistical realization. Consequently, the model cannot generalize to unseen geological structures, alternative permeability variograms, or different well configurations. 

To transition from this proof-of-concept case-specific emulator to a generalized physical solver, the training dataset must be expanded to include at least 50 to 100 simulation cases. These cases must span a wide range of geostatistical parameters (variogram ranges, anisotropy ratios, and facies distributions) and well configurations (varying injector and producer placements). A dataset of this scale would provide the spatial and geological diversity necessary for the GNN to learn generalizable transport mappings, rather than memorizing cyclic trajectories.
